% ============================================================================
% Section 6 (Comparative evaluation) -- draft for discussion
%
% Structure (the rationale is the section preamble, not a subsection):
%   6.1 What the framework costs against itself   (microbenchmark)
%   6.2 The same pipeline, five ways              (composition)
%   6.3 Keeping the output current                (freshness against cost)
%   6.4 What cannot be expressed                  (capability)
%   6.5 Threats to validity
%
% First person is kept near the paper's own density (~2 per 1000 words) and
% reserved for authorial stance: what is claimed, what is disclaimed, and the
% conflict-of-interest note in the threats.
%
% NUMBERS. Everything in 6.2 and 6.3 is now the server run (server_results/,
% 2026-09-24): experiment 1 complete (startup, rate, depth, payload, stream,
% work, work-depth, stage costs), experiment 2 as a size sweep over 5/10/20/30
% snapshots, all nine arms, three repetitions, 108 runs, every one correct.
%
% Still from the laptop, and marked in place:
%   - the interior of the granularity table (the server ran only the two ends
%     of the axis: one task for the whole input, and one task per snapshot);
%   - the whole of 6.4, freshness -- phase 3 did not run on the server.
%
% \res{} still marks a number to be replaced. Three groups remain:
%   1. the interior of the granularity table (run phase 2 on the server);
%   2. all of 6.4, freshness (run phase 3 on the server);
%   3. the CPU crossover in 6.3, which is an extrapolation past 30 snapshots --
%      extending the size sweep would turn it into a measurement.
% Plus three anecdotes carried over from laptop debugging (the Nextflow
% ordering counts and the StreamFlow gather limit), which are dataset-dependent
% rather than timings.
%
% The section is labelled sec:evaluation, which is what main.tex references.
% ============================================================================

\newcommand{\res}[1]{\textbf{[#1]}}

\section{Comparative evaluation}
\label{sec:evaluation}

The preceding sections argued that a declarative, RDF-described pipeline is
worth having. This section asks what it costs.

Table~\ref{tbl:comparison} says what each framework can express. It says nothing
about what the expressiveness costs, which is the question facing someone who
already runs a workflow system and is asked to adopt another. Two costs are
worth separating. The first belongs to the architecture: every message leaves a
processor, crosses the orchestrator and enters another runner, where a program
calling the same functions in one process would have passed a pointer. The
second belongs to the architecture in use: what it costs to run a real pipeline
this way rather than as scheduled tasks over files.

A single experiment would confound them. A comparison against another framework
cannot separate orchestration from workload, and a microbenchmark cannot say
whether the overhead matters once real processors run. The evaluation is
therefore in two parts. Section~\ref{sec:eval-micro} removes the framework and
changes nothing else. Section~\ref{sec:eval-pipeline} takes a pipeline that is
deployed and in use, and runs it under five systems. The two meet in one number:
the microbenchmark prices a message crossing the orchestrator, and the pipeline
experiment shows that this price is not what separates the systems.

Three rules hold throughout. Every system runs the same implementations: the
published processors, or command-line tools that instantiate those same classes.
No system shares warm state between invocations, since that is part of what
distinguishes them. Every run is checked against output computed from the input,
so a system that is fast because it did less work is reported as failed rather
than fast. The benchmark, the recorded input, the pipeline definitions and every
measurement behind the numbers below are published, so that a reader can rerun
them or add a system of their own.

One claim we do not make. There is no single correct way to express a streaming
workload in a batch system, and no search for an optimal encoding was attempted.
What follows instead names the parameter that separates the encodings --- how
much data one invocation handles --- and measures across its range, so that a
reader can locate their own deployment on the curve.

% ----------------------------------------------------------------------------
\subsection{What the framework costs against itself}
\label{sec:eval-micro}

The first experiment holds the processors fixed and removes the framework. Both
arms instantiate the same processor classes over the same workload. One is an
ordinary RDF-Connect deployment. The other is a driver that constructs the
processors in a single process and connects them with in-memory channels that
reproduce the runner's acknowledgement semantics, including the rule that a
write resolves only once the reader has consumed the message. The difference
between the arms is the orchestrator, the runners and the gRPC transport.

The workload is synthetic: a generator, a chain of pass-through processors and
a sink, parameterised by message count, message size, pipeline depth, and the
microseconds of CPU work each processor spends per message. The work parameter
is the interesting one. With no work per message, everything measured is
overhead; as the work grows, overhead becomes a shrinking share of a real cost.

With processors that do nothing, a message costs 0.234\,ms through RDF-Connect
and 0.002\,ms in process. The ratio is large but not the useful number, because
no processor does nothing. Adding work per message gives the shape in
Table~\ref{tab:eval-micro}: the framework adds a roughly constant
0.23--1.5\,ms per message, which is 35\,\% of the total at 1\,ms of work per
message and 7.5\,\% at 10\,ms.

\begin{table}[t]
\caption{Per-message cost with and without the framework, by work per message
(medians of three runs, depth 0).}
\label{tab:eval-micro}
\setlength{\tabcolsep}{4pt}
\begin{tabular}{rrrrr}
\toprule
work/msg & RDF-Connect & in-process & added & overhead \\
(ms) & (ms) & (ms) & (ms) & \\
\midrule
0    & 0.234 & 0.002 & 0.232 & 99\,\% \\
0.01 & 0.269 & 0.014 & 0.254 & 95\,\% \\
0.1  & 0.414 & 0.104 & 0.310 & 75\,\% \\
1    & 1.550 & 1.007 & 0.544 & 35\,\% \\
10   & 10.83 & 10.02 & 0.810 & 7.5\,\% \\
100  & 101.6 & 100.0 & 1.533 & 1.5\,\% \\
\bottomrule
\end{tabular}
\end{table}

Following Task Bench~\citep{slaughter2020}, this is summarised as a minimum
effective task granularity: the work per message at which half the time is still
framework. For RDF-Connect it is 0.47\,ms.

That threshold locates real work on the curve, which is what makes it useful.
Measured on their own, the two heaviest stages of
Section~\ref{sec:eval-pipeline} cost 0.60 and 0.39\,ms per record once warm ---
that is, on the threshold rather than safely past it. A knowledge engineering
pipeline of this kind therefore operates where the framework is neither
negligible nor dominant, and the lever that moves it is granularity: a processor
that handles a record at a time pays the overhead once per record, while one
that handles a batch of a hundred pays it once per hundred. This is why
Section~\ref{sec:eval-pipeline} measures granularity rather than assuming a
value for it.

The cost is per hop, not per pipeline. Over a chain of pass-through processors,
each additional hop adds 0.16\,ms to a message's journey, against 0.001\,ms in
process: a message crosses a depth-8 chain in 1.52\,ms where the same functions
in one process take 0.012\,ms. Depth is therefore the parameter a pipeline
author pays for, and it is the one a declarative description encourages them to
spend.

% ----------------------------------------------------------------------------
\subsection{The same pipeline, five ways}
\label{sec:eval-pipeline}

The second experiment leaves the microbenchmark's synthetic workload behind and
takes a pipeline that runs in production, published by a third party and not
written for this paper. It is expressed five times, and every
expression runs the same six processor implementations over the same
recorded input, so that the only variable is what carries a message from one
stage to the next. The pipeline, the five systems and the correctness check are described in
turn, and then what each one costs.

\begin{figure}[pos=htbp]
  \centering
  \includegraphics[width=\columnwidth]{figures/growth.pdf}
  \caption{Cost against the amount of input, for every system and every engine,
  on logarithmic axes. Colour is the system and line style is the CWL engine, so
  the three engines executing the same workflow document appear as three lines
  of one colour. Two groups separate in both panels: the systems whose cost is
  dominated by a fixed startup, which flatten as the input grows, and the
  scattered encodings, whose cost is dominated by starting a process per
  snapshot and which stay steep. In CPU time RDF-Connect starts above the batch
  workflow and the shell pipeline and converges on them; in elapsed time it does
  not, because its channels advance in lock-step.}
  \label{fig:growth}
\end{figure}

\subsubsection{The pipeline and the data}

\texttt{CA-Blue-Bike-LDES} (Section~5) publishes the availability of a Belgian
bike-sharing service as a Linked Data Event Stream. It polls an HTTP API for
the state of every station, lifts each snapshot to RDF, detects which stations
changed, annotates the changes as stream members, fragments them by time and
writes the fragments as a static LDES. Table~\ref{tab:eval-stages} lists the
stages and the components that implement them.

\begin{table}[t]
\caption{The stages of the benchmark pipeline. All six components are published
processors; five of the six KE tasks of Section~5.5 appear.}
\label{tab:eval-stages}
\begin{tabular}{llll}
\toprule
\# & Stage & Component & Runtime \\
\midrule
1 & semantic lifting & \texttt{RmlMapper} & JVM \\
2 & quality gate & \texttt{Validate} & Node.js \\
3 & change detection & \texttt{DumpsToFeed} & Node.js \\
4 & stream annotation & \texttt{Sdsify} & Node.js \\
5 & fragmentation & \texttt{Bucketize} & Node.js \\
6 & publication & \texttt{LdesDiskWriter} & Node.js \\
\bottomrule
\end{tabular}
\end{table}

The RML mapping is the deployment's own document: nine triples maps over
mobility vocabularies. One line changed, so that the mapping
reads the pipeline's channel instead of the live API. The shapes that gate the
data also describe the feed's members, which is how the deployed DCAT-AP feeds
are configured.

The input is recorded from the API the deployment consumes, so that every run
replays the same data. The unit of that recording is a \emph{snapshot}: one poll
of the API, holding the state of all 327 stations at that instant, stored as one
line of JSON. A run over $n$ snapshots replays the first $n$ lines, so thirty
snapshots is half an hour of the feed at the recorded rate of one poll a minute.
It is the quantity of input, not a measure of time, that the sweeps below vary.
Of the 327 stations, 70 carry the timestamp the mapping keys its reports on; all
70 members are created within the first five snapshots, after which the feed
grows by 2.4 updates per snapshot. That shape --- a large first load and a slow
tail --- is what a stream of this kind looks like.

\subsubsection{The systems}

\begin{description}
\item[RDF-Connect.] One pipeline description; the mapper on the JVM runner and
  the other five processors on the Node.js runner.
\item[Shell pipeline.] The six stages as command-line tools in one pipe, with
  the channels that carry stream metadata as named pipes. It runs the same
  processes as the RDF-Connect deployment with the operating system in place of
  the orchestrator, and so bounds what mediated channels and a declarative
  description can cost.
\item[CWL, batch.] One \texttt{CommandLineTool} per stage, one task each, over
  the whole input.
\item[CWL, scattered.] The same tools, with the input split into chunks and the
  stateless stages run once per chunk. CWL does not split files, so the split
  and the concatenation are steps of the workflow and are measured with it.
\item[Nextflow.] The same tools again, composed with channels. A process still
  runs once per item, but consumers start as items arrive, so the mapping of one
  snapshot overlaps the validation of the previous one.
\end{description}

The command-line tools are not reimplementations. A small adapter instantiates a
published processor from an RDF description of the step, materialises its
arguments through the processor's own SHACL shapes exactly as the runner does,
and connects its channels to standard input and output. Every system therefore
runs the same processor code with the same configuration; what differs is what
carries messages between stages.

Both CWL encodings run on three engines.\footnote{\texttt{cwltool}
3.2.20260720092025, Toil 9.5.0 and StreamFlow 0.2.0rc3, from the same workflow
documents, without containers. Figure~\ref{fig:growth} plots all three, and per-engine
figures are in the artefact.} Unless stated otherwise, the reported CWL figures
are cwltool's.

\subsubsection{Correctness}

Each run is checked against output derived from the input rather than from
another system. A member is created the first time a station is observed with a
given timestamp, and updated when the data the member shape reaches changes;
both are computable from the recorded snapshots. What the shape reaches matters:
the station's total capacity, two properties away from the member, changes
without producing an update, and a check written against the station records
rather than the shape reports failures that are not there. Timestamps the
pipeline mints itself are excluded, since they differ between runs by design.
All 108 runs reported here published exactly the members and activities the
input implies.

\subsubsection{Setup}

One machine: a four-vCPU KVM virtual machine (QEMU virtual CPU, single socket,
no simultaneous multithreading, 16\,MiB L3), 15\,GiB RAM, Ubuntu 24.04.4 on
Linux 6.8.0, with Node.js 24.21.0, OpenJDK 21.0.12.1 and Nextflow 26.04.6. Four cores is
few enough to matter, and is returned to below. Working directories are on tmpfs,
so the file-passing systems are not measured against the disk.
Each configuration runs three times; medians are reported. Repetitions spread by
at most 16\,\% of CPU time, at the smallest input where a run is mostly process
startup, and by under 3\,\% for the scattered encodings, where the per-task cost
dominates.

Each run executes inside a transient cgroup, from which CPU time and peak memory
are read. Sampling would not do here: a scattered workflow starts thousands of
processes that live for well under a second, and a process that begins and ends
between two samples contributes nothing to the total. Intermediate storage is
the largest total size of a system's working directories during the run, sampled
every 250\,ms.

\subsubsection{Results}

\begin{table}[t]
\caption{The pipeline over 30 snapshots. Elapsed is wall-clock time from launch
to last output, the quantity a scheduling paper would call makespan; \emph{per snap.} is the marginal CPU cost of one further
snapshot, fitted over the sweep of Figure~\ref{fig:growth}. Mean cores busy is
given in the text rather than here. The scattered
encoding is at its finest setting, one snapshot per task, and the CWL rows are
cwltool's. Medians of three runs; every run published the expected members and
activities.}
\label{tab:eval-systems}
\setlength{\tabcolsep}{4pt}
\begin{tabular}{lrrrrr}
\toprule
System & elapsed & CPU & per snap. & mem & disk \\
 & (s) & (s) & (s) & (MB) & (MB) \\
\midrule
Shell pipeline & 16.0  & 51.2  & 0.92  & 1645 & --- \\
RDF-Connect    & 35.5  & 62.4  & 0.78  & 1707 & --- \\
\midrule
CWL, batch     & 31.9  & 51.9  & 0.87  & 823  & 21.4 \\
CWL, scattered & 127.9 & 402.0 & 13.04 & 673  & 32.8 \\
Nextflow       & 145.2 & 441.4 & 13.34 & 998  & 47.7 \\
\bottomrule
\end{tabular}
\end{table}

Table~\ref{tab:eval-systems} reports every system at the largest input, and
Figure~\ref{fig:growth} shows how each one got there.

At 30 snapshots RDF-Connect costs 20\,\% more CPU than the batch workflow and
22\,\% more than the shell pipeline, and a sixth of what the same pipeline costs
when CWL is made to work a snapshot at a time.

The interesting quantity is not the total but how it grows, which is what
Figure~\ref{fig:growth} makes legible. CPU time is close to linear in the input
for every system, so fitting a fixed and a per-snapshot cost separates two
effects that the totals confound. RDF-Connect has the \emph{lowest} marginal
cost of any system measured, 0.78\,s per snapshot against 0.87 for the batch
workflow and 0.92 for the shell pipeline; what it pays for is a large fixed
cost, 41.1\,s against 27.2 and 24.3, being two runners and a JVM that the
file-passing systems also start but that the shell pipeline starts more cheaply.
Its disadvantage therefore narrows as the input grows --- from 41\,\% more CPU
than the batch workflow at five snapshots to 20\,\% at thirty, the convergence of
the blue, green and grey lines in the upper panel --- and the fits cross at
\res{roughly 150 snapshots}, about two and a half hours of this feed. That
crossing lies beyond the measured range and is flagged as such; the sweep should
be extended past it.

The scattered encodings never flatten, because their cost is charged per task
rather than per run. Their marginal cost is 13.0 and 13.3\,s per snapshot,
roughly fifteen times the streaming systems', and their lines stay steep across
the sweep.

The three CWL engines appear as three lines of one colour, which is the point of
plotting them together. cwltool and StreamFlow are indistinguishable at this
scale. Toil, the dotted line, tracks the same shape at a constant offset: it
spends 1.3--1.5$\times$ their CPU over 1.8--2.5$\times$ their elapsed time while
running fewer cores, so it is slower and more expensive than the other two
without being a different kind of system. A difference between engines is thus a
property of the engine, not of the description they share.

Elapsed time and CPU time diverge. The shell pipeline finishes in less than half
RDF-Connect's time while running 3.2 cores to its 1.8. What separates the two is
concurrency rather than transport: a message crosses the orchestrator in
0.23\,ms (\S\ref{sec:eval-micro}), which over 30 snapshots cannot account for
20\,s, but a channel admits one message at a time, so the stages advance in
lock-step where a shell pipe lets them run ahead of each other into kernel
buffers. The cost of the architecture is paid in latency and in startup, not in
the transport itself.

Four cores is enough to divide the systems in two. Dividing a run's CPU time by
the core count gives the elapsed time it would have if it saturated the machine;
the shell pipeline comes within 25\,\% of that floor, the scattered encoding
within 27\,\% and Nextflow within 32\,\%, so all three are CPU-bound here and
would finish sooner on a larger machine without spending any less CPU.
RDF-Connect and the batch workflow sit at 2.3 and 2.5$\times$ the floor: both
leave half the machine idle, one because its channels advance in lock-step and
the other because it runs one stage at a time. On this axis RDF-Connect is not
an outlier among the declarative systems --- it is the file-passing batch
workflow's equal --- and the shell pipeline is the exception.

Storage grows with the input for every system that passes files, and stays at
zero for the streaming ones: 4.1 to 21.4\,MB for the batch workflow over this
range and 8.6 to 47.7\,MB for Nextflow, which keeps a working directory per
task. Each writes every stage's output, and here also every stage's state, where
a streaming system passes both in memory. Peak memory runs the other way: the
streaming systems hold the whole pipeline resident at once (1.6--1.7\,GB), the
file-passing ones hold one stage at a time (0.7--1.0\,GB).


\paragraph{Granularity.} How much data one invocation handles is the parameter
that distinguishes the encodings, and the one a workflow author sets by choosing
a chunk size. Table~\ref{tab:eval-unit} sweeps it, from the whole input in one
task to one snapshot per task. Across the sweep the work is identical and the
output is the same, yet CWL spends 7.7$\times$ the CPU at the fine end and
Nextflow 8.5$\times$, because every task starts its own JVM and Node.js runtime.
The microbenchmark prices that directly: one extra start of the mapper costs as
much as 1\,700 warm records through it, and of the validator 700, so a chunk
size that adds a process start per snapshot buys nothing and pays for two. The
penalty also grows with the input, since it is charged per task rather than per
run: 2.6$\times$ at five snapshots, 7.7$\times$ at thirty. RDF-Connect works at
one snapshot per message throughout, for 1.20$\times$ the CPU of the batch
workflow and a sixth of what CWL spends to match its granularity.

\begin{table}[t]
\caption{Granularity: CPU time and elapsed time (s) over 30 snapshots, by snapshots
per task; the last column is the whole input in one task. RDF-Connect and the
shell pipeline work a snapshot at a time throughout and are given for reference. \res{The interior columns are from the
earlier laptop run; the server measured only the two ends of the axis.}}
\label{tab:eval-unit}
\begin{tabular}{lrrrr}
\toprule
snapshots per task & 1 & \res{5} & \res{10} & 30 \\
\midrule
CWL, CPU (s)       & 402.0 & \res{--} & \res{--} & 51.9 \\
CWL, elapsed (s)   & 127.9 & \res{--} & \res{--} & 31.9 \\
Nextflow, CPU (s)  & 441.4 & \res{--} & \res{--} & \res{--} \\
Nextflow, elapsed (s) & 145.2 & \res{--} & \res{--} & \res{--} \\
\midrule
RDF-Connect    & \multicolumn{4}{c}{62.4 CPU, 35.5 elapsed} \\
Shell pipeline & \multicolumn{4}{c}{51.2 CPU, 16.0 elapsed} \\
\bottomrule
\end{tabular}
\end{table}

Nextflow's channels let consumers start before their producers finish, which
shortens its elapsed time relative to its own CPU time, but overlap does not recover
the cost of a process per item: it spends 10--37\,\% more CPU than the scattered
CWL workflow, the gap narrowing as the tasks grow. Its advantage would grow with
the work per task, and these tasks are seconds long.

\paragraph{The descriptions.} Counting lines that are neither blank nor comment,
describing this pipeline takes 105 lines of RDF in a single document for
RDF-Connect and 206 lines across seven files for CWL, where every stage needs a
tool description besides its configuration and every stateful stage needs its
state wired between tasks. Inserting a further processor, skolemisation, changes
10 lines of the one document against 31 lines across three files: seven in the
workflow, a new fifteen-line tool description, and a nine-line description of
the step's configuration that the command-line arms need and the RDF-Connect
pipeline carries inline. Both modified definitions were executed rather than
only counted.

What that comparison measures needs stating, because it is not total RDF against
total YAML. The RDF-Connect pipeline additionally imports 516 lines of processor
descriptions --- the SHACL contract and metadata each processor package ships,
plus two more that declare themselves inside their jars --- and those lines are
not in the 105. They are written once by the processor author and reused by
every pipeline that instantiates the processor, whereas CWL's tool descriptions
sit in the 206 because they were written for this pipeline. The honest reading
is therefore narrower than a ratio of file sizes: for a pipeline assembled from
processors that have already been described, an author writes about half as
much, and the saving comes from the ecosystem of Section~5 rather than from RDF
being terser than YAML. Where no published description exists, someone writes
one, and the 516 lines are what that costs.


% ----------------------------------------------------------------------------
\subsection{Keeping the output current}
\label{sec:eval-fresh}
% NOT YET RUN ON THE SERVER. Every number below is the laptop run of 2026-09-23
% (results/freshness.json); phase 3 of bench/all-bluebike.sh has to be repeated
% on the benchmark machine before this subsection is final. The shell arm's
% clock also starts at process spawn, which charges JVM startup to its first
% snapshots.

The pipeline of Section~\ref{sec:eval-pipeline} exists to keep a published
stream current. Its deployment polls every fifteen seconds and publishes
continuously. A batch system cannot do that: a step begins when its inputs exist
as complete files, and a source that never ends has no complete input file. What
it can do is run again, and the usual arrangement is a scheduled job that
rebuilds from the data accumulated so far --- the idiom of a cron entry or a
scheduled CI workflow, which needs no state to travel between runs and leaves
every invocation self-contained.

That arrangement is measured against the streaming one. Snapshots arrive at a
fixed rate; the streaming systems consume them as they arrive, and the batch
workflow is re-run every $k$ arrivals over the whole history. Freshness is the
time from a snapshot's arrival to the publication of the activities it causes;
cost is the CPU spent over the window.

\begin{table}[t]
\caption{Mean freshness and CPU cost over \res{12} snapshots arriving every
\res{3}\,s; the rebuild rows re-run every $k$ arrivals over the whole history. \res{Laptop run;
arrivals are accelerated relative to the deployment's fifteen-second poll, so
the shape transfers and the absolute latency does not.}}
\label{tab:eval-fresh}
\setlength{\tabcolsep}{4pt}
\begin{tabular}{lrrr}
\toprule
& runs & freshness (s) & CPU (s) \\
\midrule
RDF-Connect, streaming & 1 & \res{0.3}  & \res{17.5} \\
Shell pipe, streaming  & 1 & \res{0.4}  & \res{15.0} \\
Rebuild every 2        & 6 & \res{10.7} & \res{86.5} \\
Rebuild every 4        & 3 & \res{11.5} & \res{38.7} \\
Rebuild every 6        & 2 & \res{14.0} & \res{23.6} \\
\bottomrule
\end{tabular}
\end{table}

Two things follow. A scheduled rebuild cannot be fresher than its own run takes,
so its freshness has a floor that streaming does not have: RDF-Connect publishes
\res{0.3}\,s after a snapshot arrives, while spending less CPU than any of the
schedules that come closest to it. And the rebuild's cost grows with the
history: running every $k$ arrivals reprocesses about $n^2/2k$ snapshots where
streaming processes $n$, so the penalty widens as the day goes on ---
\res{18\,000 snapshot-processings against 1\,440 over a day of this feed with
hourly runs}. A deployment that resets its history daily bounds that growth,
which is a reasonable thing to do and worth saying: the rebuild idiom is
workable exactly when history is bounded by a cycle of that kind.

The alternative is to carry the state between runs and process only what is new.
Measured here, it removes the quadratic term at \res{11--44}\,\% less CPU,
and in exchange the state of three stages must travel between
invocations as files, which is the arrangement
Section~\ref{sec:eval-capability} describes. A rebuild also re-mints the
publication timestamps of members it has already published, so consumers of an
append-only stream see the whole feed change.

% ----------------------------------------------------------------------------
\subsection{What cannot be expressed}
\label{sec:eval-capability}

Some differences are not a matter of cost, because one system cannot express the
workload at all. They are reported as boundaries.

\paragraph{State across tasks.} Three of the six stages keep state: change
detection remembers the previous version of every member, the fragmenter
remembers its buckets, and the publisher appends to a tree it has already
written. A streaming processor holds that state for the life of the pipeline. A
scattered task, by contrast, is an isolated execution over staged inputs, and
neither CWL nor Nextflow can pass a value from one such task to the next. These
three stages therefore run as single tasks over the gathered stream in both
systems, and their state travels between stages as staged directories and files.
Only the two stateless stages at the front of the pipeline can be distributed,
and the state a streaming processor holds in memory appears in the other systems
as intermediate storage. This follows from the execution model, so it holds for
workflow managers of this kind generally.

\paragraph{Order.} Change detection compares each snapshot against the one
before it, which presumes the stream arrives in order. A shell pipe and an
RDF-Connect channel preserve order; a Nextflow channel does not, and gathering
its results as they complete produced \res{17} updates where \res{12} had
occurred. The workflow is correct once the gather is sorted explicitly, but the
failure is silent: the output remains well-formed RDF and the member count is
unchanged.

\paragraph{Portability of a valid description.} The same valid CWL workflow did
not behave identically on three engines. StreamFlow wrote a tool's standard
error into the file capturing its standard output, and staged a gathered array
of files into one directory, where files with the same name overwrote each
other; both silently produced wrong output. It also passes a step's command line
to \texttt{sh} as a single argument, so a gather over \res{1\,000} files exceeds
the operating system's limit and fails after every task has run. Toil retries a
failed job with more memory by default, which hides failures in timings. The
workarounds were small, and harmless on the other engines, but a portable
description did not give us portable execution.

% ----------------------------------------------------------------------------
\subsection{Threats to validity}
\label{sec:eval-threats}

\emph{One implementation is not a paradigm.} The identical workflow ran on three
CWL engines, two of which agree to within 3\,\% of CPU time on every
configuration, and the same pipeline ran on Nextflow, whose execution model
differs. The costs reported here are therefore properties of
executing a pipeline as scheduled tasks over files, not of any one engine.

\emph{We wrote RDF-Connect.} The parameter grid was fixed before the final runs;
the benchmark, the recorded input and all raw measurements are published; and
the cases where RDF-Connect loses are reported. It spends more CPU than both the
shell pipeline and the batch workflow at every size measured, and more than twice the
elapsed time of the shell pipeline; the claim drawn from the growth fit is about
the trend, not about the totals.

\emph{The inputs are small relative to the fixed costs.} At thirty snapshots the
fitted fixed cost is still the larger part of every streaming system's total, so
the sweep measures amortisation as much as throughput. The crossover reported above lies beyond the measured range and
should be read as a projection.

\emph{Parallelism differs, on a small machine.} The systems take different
shares of the four cores available: 3.2 for the shell pipeline, 3.1 for the
scattered workflow, 1.8 for RDF-Connect. Makespan alone would flatter whichever
takes more, so CPU time from the
kernel's cgroup accounting is reported alongside it, covering every descendant process including the short-lived ones a scattered
workflow starts. Four cores also caps what the scattered encodings can show:
they saturate the machine, so their elapsed
time here is close to their CPU time divided by the core count, and a wider
machine would narrow the elapsed-time gap while leaving the CPU gap --- which is what granularity actually buys ---
untouched.

\emph{Encoding choice.} There is no canonical way to run a streaming workload in
a batch system. What is reported is a family of encodings along one
parameter rather than a single point, and the granularity sweep shows what the choice is worth.

\emph{One machine, one dataset.} The pipeline is one deployment among the
several of Section~5, and the recorded input covers one API polled once a
minute. Its stations change slowly --- all 70 members are created in the first
five snapshots and the feed then grows by 2.4 updates per snapshot --- so it is
dominated by the first load; a faster-moving source would put more load on the
stages after change detection. The microbenchmark of
Section~\ref{sec:eval-micro} is what generalises the per-message cost beyond it.
